\newcommand\version{v1}
\problemname{Cookies}
Ann Britt-Caroline has $N$ different types of cookies. She has $A_i$ cookies of type $i$. Now, Ann is wondering how many cookies she has in total.

\section*{Example}
In this example, we have $N = 3$ different types of cookies. The number of cookies of each type is $3, 1, 5$.

\begin{figure}[h!]
  \centering
  \includegraphics[width=0.8\textwidth]{sample.png}
  \caption{We have three cookies of the first kind, one of the second kind, and five of the third}
\end{figure}

As pictured, this gives a total of $9$ different cookies.

\section*{Task}
You will be given $N$ and $A$. Write a program to help Ann compute the total number of cookies. You should implement the function \texttt{cookies(N, A)}:
\begin{itemize}
  \item \texttt{cookies(N, A)} - this function will be called exactly once by the judge.
  \begin{itemize}
    \item \texttt{N}: the number of types of cookies.
    \item \texttt{A}: a vector of length $N$. $A[i] (0 \le i < N)$ contains the number of cookies of type $i$. $A[i]$ is always positive.
    \item The function should return the number of cookies Ann has.
  \end{itemize}
\end{itemize}

\section*{Subtasks}
The problem consists of a number of subtasks. Each subtask gives some amount of points, and to pass
the subtask you must pass all the test cases in the subtask.

\begin{tabular}{|l|l|l|}
  \hline
  \textbf{Subtask} & \textbf{Points} & \textbf{Limits} \\ \hline
  1 & 27 & $1 \le N \le 1\,000$,  $A[i] = 1$. \\ \hline
  2 & 34 & $1 \le N \le 1\,000$,  $A[i] \le 100\,000$. \\ \hline
  3 & 39 & $1 \le N \le 100\,000$,  $A[i] \le 100\,000$. \\ \hline
\end{tabular}

\section*{Input format}
The sample judge reads input in the following format:

\begin{itemize}
  \item line 1: \texttt{N}
  \item line 2: \texttt{A[0] A[1] ... A[N - 1]}
\end{itemize}

\section*{Output format}
The sample judge writes the return value of \texttt{cookies(A, N)}.
